%% 01_introduccion.tex — Context, challenges, related work, contributions.

\section{Introduction}
\label{sec:intro}

Dynamic Application Security Testing (DAST) evaluates a running system by
sending manipulated requests and observing the response, with no access to
source code. Compared to static analysis, its main advantage is the absence
of assumptions about the target's language or framework; its main drawback is
cost. Comparative studies have repeatedly shown that black-box scanners
exhibit uneven coverage and elevated false-positive
rates~\cite{doupe2010johnny,bau2010state,parvez2015analysis}, which erodes the
development team's trust and, in practice, leads reports to be ignored.

We identify three concrete challenges that this work addresses explicitly:

\begin{enumerate}
  \item \textbf{Accumulated latency.} A realistic scan involves on the order
  of $10^4$--$10^5$ HTTP requests. Executed sequentially, network latency
  dominates total run time and makes integrating DAST into every \emph{commit}
  infeasible.
  \item \textbf{Noise (false positives).} Detection based solely on a
  reflected string or a database error pattern produces findings that are not
  exploitable. Telling an inert reflection apart from real execution requires
  a second, targeted observation.
  \item \textbf{Combinatorial explosion.} The product (injection points)
  $\times$ (payloads per category) $\times$ (evasion transformations) grows
  until testing everything under a fixed time budget becomes impossible; the
  order in which that budget is spent determines what gets detected. This
  paper's architecture is designed to withstand that growth (\cref{sec:pipeline}),
  but the evaluation in \cref{sec:resultados} exercises the scheduling
  mechanics on a constrained corpus --- \num{338}~\code{GET} parameters that
  saturate at \num{20}~requests each --- not a payload space large enough to
  demonstrate the explosion itself; \cref{sec:amenazas} discusses this
  restriction.
\end{enumerate}

\subsection{Related work}
\label{sec:related}

\noindent\textbf{LLM-based vulnerability triage.} LLM-assisted triage of
security findings is an active area outside of DAST. Static-analysis
pipelines that pair a deterministic scanner (Semgrep, CodeQL) with an LLM
adjudication stage report large false-positive reductions: Munson~\emph{et
al.} enrich static findings with reachability context before LLM
adjudication~\cite{munson2025llmfriends}; Ameen~\emph{et al.}'s multi-agent
\textsc{QASecClaw} raises the $F_1$ score of standalone Semgrep from
78.4\,\% to 90.9\,\% on Java test cases by suppressing non-exploitable
alerts~\cite{ameen2026qasecclaw}; Agrawal and Ahi's \textsc{SAST-Genius}
hybrid framework cuts Semgrep's false positives by 91\,\% (225 to 20) through
LLM-driven triage, dynamic bug descriptions, and exploit-generation
validation~\cite{agrawal2025sastgenius}; and Du~\emph{et al.}'s industrial
study at Tencent reports that hybrid LLM/static-analysis triage eliminates
94--98\,\% of false positives while preserving recall, at a marginal cost of
\SIrange{2.1}{109.5}{\second} per alert~\cite{du2026reducing}. These results
motivate the LLM-triage stage evaluated in \cref{sec:ai-triage}: unlike the
static-analysis setting, a DAST finding already carries live request/response
evidence, so triage here re-scores an executed probe rather than a static
alert. All four studies report their false-positive reductions against
corpora with a substantial false-positive population to begin with; none
addresses the scheduling problem that determines \emph{which} payloads reach
the target under a fixed request budget in the first place, nor evaluates
triage on a target whose noise floor is already low --- both of which
\cref{sec:resultados,sec:ai-triage} take up directly.

\noindent\textbf{Payload and seed scheduling.} Outside DAST, the greybox
fuzzing literature treats seed/input scheduling as a first-class problem:
She~\emph{et al.} prioritize fuzzing seeds by graph-centrality analysis of
the program's control-flow graph to maximize coverage growth per
execution~\cite{she2022effective}, the same anticipation argument this paper
makes for confidence-based payload prioritization
(\cref{sec:pipeline,sec:res-sintesis}) --- earlier detection per request
under a fixed budget, not necessarily a higher final yield. To our knowledge
no prior work transfers this scheduling framing to black-box DAST payload
selection, where the search space is a fixed, tester-supplied corpus rather
than a mutation-generated one.

\noindent\textbf{OWASP Benchmark as an evaluation target.} Potti~\emph{et
al.} benchmark OWASP ZAP v2.12.0 and v2.13.0 on the OWASP Benchmark project,
reporting an aggregate true-positive rate of 68.75--74.63\,\% at a
0.86--3.45\,\% false-positive rate across vulnerability
categories~\cite{potti2025benchmarking}. \Cref{sec:res-external} runs a newer
ZAP release (2.17.0) ourselves against this paper's own \num{338}-target
list under the same oracle used throughout this paper, and finds their
published range consistent with our own single-run result.

\noindent\textbf{Contributions.} This paper makes the following
contributions:

\begin{itemize}
  \item An asynchronous DAST engine with a single HTTP choke point,
  back-pressured concurrency control, token-bucket rate limiting, and
  safeguards against SSRF in redirect following (\cref{sec:arquitectura}).
  \item A payload-scheduling pipeline with decoupled stages --- context
  interleaving, prioritization, ordering, and a live adaptive-sorting
  feedback loop --- plus a runtime-confirmation stage that calibrates finding
  confidence (\cref{sec:arquitectura,sec:metodologia}).
  \item An optional LLM-based triage stage that re-scores executed probes
  using request/response context, and an honest evaluation of it against the
  same budget grid, including a null-result case where it suppresses nothing
  and a measured wall-clock overhead of zero (\cref{sec:ai-triage}).
  \item A non-blocking \code{TelemetryWorker} that logs one JSONL row per
  payload probe, enabling instance-by-instance detection-cost measurement
  (\cref{sec:telemetria}).
  \item A reproducible evaluation methodology: a containerized testbed, an
  independent grading oracle, an RCBD experimental design, and a
  reproducibility manifest with a corpus hash (\cref{sec:metodologia}).
  \item An ablation study that isolates each pipeline stage's contribution to
  the detection-versus-cost curve and to the false-positive rate
  (\cref{sec:resultados}).
  \item An external baseline against OWASP ZAP on the identical target list
  under the same oracle, reported honestly including the categories where
  ZAP outperforms \herramienta (\cref{sec:res-external}).
\end{itemize}

The rest of the paper is organized as follows. \Cref{sec:arquitectura}
describes the engine's design. \Cref{sec:metodologia} details the testbed and
the experimental design. \Cref{sec:resultados} presents the ablation
evaluation and \cref{sec:ai-triage} the LLM-triage evaluation.
\Cref{sec:amenazas} discusses threats to validity, and \cref{sec:conclusion}
concludes.
